\documentclass[12pt]{article}
\usepackage[utf8]{inputenc}
\usepackage[T1]{fontenc}
\usepackage{hyperref}
\title{Texmelucan Zapotec Language Resources for K-12 Curriculum Development}
\author{Compiled by Claude Sonnet 5.5 (claude-sonnet-5-5) for Andr\'as Kornai}
\date{}
\begin{document}
\maketitle

\section{Introduction}
Texmelucan Zapotec (ISO 639-3 zpz, Glottocode texm1235) is an Otomanguean language spoken in San Lorenzo Texmelucan, Oaxaca, Mexico. It is also called Central Sola de Vega Zapotec and Papabuco. Estimates of the number of speakers are about 4,100 to 4,600. Most of what exists is linguistic research and religious material. Nothing aimed at schools turned up.

\section{Primary Educational Resources}
No primers, textbooks or school programmes were found in search results.

\section{Linguistic and Reference Materials}
\begin{itemize}
\item Wikipedia overview, including a summary of the phonology: \url{https://en.wikipedia.org/wiki/Texmelucan_Zapotec}
\item Glottolog entry: \url{https://glottolog.org/resource/languoid/id/texm1235}
\item Charles H. Speck, ``The Phonology of the Texmelucan Zapotec Verb'', \emph{International Journal of American Linguistics} (1984), cited in the Wikipedia article. A separate book, \emph{Zapotec oral literature; El folklore de San Lorenzo Texmelucan} (SIL, 1998), is also cited there. No direct links to either were found.
\end{itemize}

\section{Religious and Community Resources}
\begin{itemize}
\item Scripture Earth: Bible texts in Texmelucan Zapotec, including a New Testament PDF: \url{https://www.scriptureearth.org/00eng.php?iso=zpz}
\end{itemize}

\section{Digital and Audio Resources}
\begin{itemize}
\item Global Recordings Network, Zapotec, Texmelucan: evangelism and Bible-teaching audio: \url{https://www.globalrecordings.net/en/language/zpz}
\end{itemize}

\section{Curriculum Development Context}
The language is thought to be spoken by the whole community, but there is no direct evidence of this. It is not known to be taught in schools. The New Testament and Speck's folklore collection could supply native-language text. The phonological descriptions would help in designing an orthography. No dictionary or learner course was found, so those would still need to be developed.

\end{document}